\begin{statementbox}
	\textbf{\textcolor{CStatement}{条件}}
	\begin{description}[style=nextline, leftmargin=2.8em, labelsep=0.5em, itemsep=0.35em]
		\item[\textbf{\textcolor{CStatement}{条件 1（向量法元素）：}}] 平面 $\alpha$ 过点 $A$，$\mathbf{n}$ 是平面 $\alpha$ 的一个非零法向量。
	\end{description}

	\textbf{\textcolor{CStatement}{结论}}
	\begin{description}[style=nextline, leftmargin=2.8em, labelsep=0.5em, itemsep=0.45em]
		\item[\textbf{\textcolor{CStatement}{结论一（距离公式）：}}]
			\textbf{\textcolor{CStatement}{直观描述：}}
			点 $P$ 到平面 $\alpha$ 的距离等于向量 $\vec{AP}$ 在法向量方向投影的绝对值。
			\[
				d = \frac{|\vec{AP}\cdot\mathbf{n}|}{|\mathbf{n}|}.
			\]

		\item[\textbf{\textcolor{CStatement}{推广结论（平面方程形式）：}}]
			\textbf{\textcolor{CStatement}{直观描述：}}
			若平面用一般方程表示，则代入点坐标绝对值除以系数模长即得距离。
			\[
				d = \frac{|Ax_P+By_P+Cz_P+D|}{\sqrt{A^2+B^2+C^2}},
			\]
			其中平面 $\alpha$ 的方程为 $Ax+By+Cz+D=0$。
	\end{description}
\end{statementbox}
